\documentclass[11pt]{article}
\usepackage[margin=1in]{geometry}
\usepackage{amsmath,amssymb,amsthm}
\usepackage{booktabs}
\usepackage{hyperref}
\newtheorem{theorem}{Theorem}
\newtheorem{lemma}[theorem]{Lemma}
\newtheorem{proposition}[theorem]{Proposition}
\newtheorem{corollary}[theorem]{Corollary}
\theoremstyle{definition}
\newtheorem{definition}[theorem]{Definition}
\theoremstyle{remark}
\newtheorem{remark}[theorem]{Remark}
\newcommand{\et}{\tilde e}
\newcommand{\ft}{\tilde f}
\newcommand{\wt}{\operatorname{wt}}

\title{Disproof of Conjecture 00000003824:\\
the Nakajima monomial crystal $M(\lambda)$ is not $\bigoplus_{\mu\le\lambda}B(\mu)$}
\author{jilint777}
\date{2026-10-04}

\begin{document}
\sloppy
\maketitle

\begin{abstract}
Conjecture 00000003824 asserts that the Nakajima monomial crystal $M(\lambda)$
decomposes into connected components $\bigoplus B(\mu)$, where $\mu$ runs exactly
over the dominant integral weights $\mu\le\lambda$ (dominance order), each
appearing exactly once. This is false already for $\mathfrak{sl}_2$ and
$\lambda=2$: the monomial crystal generated by the highest-weight monomial
$Y_0^2$ is $\{Y_0^2,\ Y_0Y_1^{-1},\ Y_1^{-2}\}\cong B(2)$, a single string with
$3$ vertices, while $B(2)\oplus B(0)$ has $4$ vertices and a highest-weight
vertex of weight $0$. More generally $M(\lambda)\cong B(\lambda)$ for every
$\lambda$, so the claim fails for every $\lambda\ge2$. We also treat two
alternative readings of ``$M(\lambda)$'' (the union of the components of all
highest-weight monomials of weight $\lambda$, and the crystal of products of
$\lambda$ fundamental monomials); $\lambda=3$ defeats all three readings
simultaneously. The $\mathfrak{sl}_2$ counterexample, the failure for all
$\lambda\ge2$, and the three readings at $\lambda=3$ are formalized in Lean~4
(core library only).
\end{abstract}

\section{The conjecture and the clause we refute}
\begin{quote}
\emph{Definition:} The Nakajima monomial crystal $M(\lambda)$ denotes the
highest-weight monomial graph generated by monomial variables under Chevalley
rules. \emph{Conjecture:} $M(\lambda)$ decomposes into connected components
$\bigoplus B(\mu)$, where $\mu$ runs exactly over dominant integral weights
bounded by $\lambda$ in the dominance order, each appearing exactly once.
\end{quote}
The Chinese version says the same (``the connected-component decomposition of
$M(\lambda)$ is $\bigoplus B(\mu)$, $\mu$ ranging exactly over the dominant
integral weights not exceeding $\lambda$ in dominance order, each exactly
once''). No Lie algebra is fixed, so the statement is universal over
(symmetrizable) Kac--Moody algebras and dominant $\lambda$; a single
algebra and a single $\lambda$ suffice for a disproof. We take
$\mathfrak g=\mathfrak{sl}_2$.

For $\mathfrak{sl}_2$ the dominant integral weights are $\mu=m\varpi$ with
$m\in\mathbb N$; we write simply $\mu=m$. The simple root is $\alpha=2$ and
$\mu\le\lambda$ in dominance order means $\lambda-\mu\in\mathbb N\alpha$, i.e.
$\mu\le\lambda$ and $\lambda-\mu$ even. Hence the conjecture predicts
\[
M(\lambda)\ \cong\ B(\lambda)\oplus B(\lambda-2)\oplus B(\lambda-4)\oplus\cdots,
\qquad\text{in particular}\qquad M(2)\cong B(2)\oplus B(0).
\]
\textbf{Clause refuted:} for $\mathfrak g=\mathfrak{sl}_2$ and $\lambda=2$,
$M(2)\not\cong B(2)\oplus B(0)$ (Theorem~\ref{thm:main}); in fact
$M(\lambda)\not\cong\bigoplus_{\mu}B(\mu)$ for every $\lambda\ge2$
(Theorem~\ref{thm:all}).
If ``bounded by $\lambda$'' is read as the plain inequality $\mu\le\lambda$
without the parity condition ($\mu\in\{0,1,2\}$ for $\lambda=2$), the claim
still fails, and even more so: $M(\lambda)\cong B(\lambda)$ (Theorem~\ref{thm:all})
misses $B(\lambda-1)$ for every $\lambda\ge1$ and also $B(\lambda-2)$ for $\lambda\ge2$
(Lean: \texttt{main\_fails\_all'}).

\section{Definitions}
\subsection{The monomial crystal in type \texorpdfstring{$A_1$}{A1}}
We use Kashiwara's version~\cite{kas} of Nakajima's monomial
crystal~\cite{nak}; see also~\cite{kks}. In rank one the constants $c_{ij}$
of~\cite{kas} play no role, and Nakajima's variables $Y_{aq^{k}}$ are identified
with Kashiwara's $Y_m$ by re-indexing (Remark~\ref{rem:conv}).

Let $\mathcal M$ be the set of Laurent monomials $M=\prod_{m\in\mathbb Z}Y_m^{y_m}$
($y_m\in\mathbb Z$, finitely many nonzero), and $A_m=Y_mY_{m+1}$. Put
\[
\wt(M)=\sum_m y_m,\qquad
\varphi(M)=\max_{n\in\mathbb Z}\sum_{k\le n}y_k,\qquad
\varepsilon(M)=\max_{n\in\mathbb Z}\Bigl(-\sum_{k>n}y_k\Bigr).
\]
Both maxima are $\ge0$ (take $n$ below, resp.\ above, all indices) and
$\varphi(M)-\varepsilon(M)=\wt(M)$. The Kashiwara operators are
\[
\ft M=\begin{cases}A_{n_f}^{-1}M,& \varphi(M)>0,\\ 0,&\varphi(M)=0,\end{cases}
\qquad
\et M=\begin{cases}A_{n_e}M,& \varepsilon(M)>0,\\ 0,&\varepsilon(M)=0,\end{cases}
\]
where $n_f=\min\{n:\sum_{k\le n}y_k=\varphi(M)\}$ and
$n_e=\max\{n:-\sum_{k>n}y_k=\varepsilon(M)\}$. (When $\varphi(M)=0$ the set
defining $n_f$ is unbounded below, so the formula gives no $n_f$; the value
$\ft M=0$ is the one under which $\varphi$ is the length of the $\ft$-string,
and it is the standard semi-normal convention (see Kang--Kim--Shin~\cite{kks};
Kashiwara~\cite{kas} shows that $\mathcal M$ is semi-normal). The same remark
applies to $\et$.) $\mathcal M$ is a crystal~\cite{kas}; in particular
$\ft M=M'\iff \et M'=M$.

\begin{definition}
For a dominant weight $\lambda$, $M(\lambda)$ is the subcrystal of $\mathcal M$
generated by the highest-weight monomial $Y_0^{\lambda}$: the set of monomials
obtained from $Y_0^\lambda$ by finite sequences of $\et$ and $\ft$ (equivalently,
since $\et$ and $\ft$ are mutually inverse, the connected component of
$Y_0^\lambda$ in the crystal graph). This is the standard meaning of the
notation $M(\lambda)$ in the literature~\cite{kas,kks}, and it is what the
conjecture's definition describes (``the highest-weight monomial graph generated
\dots under Chevalley rules''). We call it the \emph{main reading}.
\end{definition}

\subsection{The target crystal}
$B(\mu)$ ($\mu\in\mathbb N$) is the crystal of the irreducible
$\mathfrak{sl}_2$-module of highest weight $\mu$: the string
$b_0\to b_1\to\cdots\to b_\mu$ with $\wt b_k=\mu-2k$, $\ft b_k=b_{k+1}$ ($k<\mu$),
$\ft b_\mu=0$, $\et b_k=b_{k-1}$ ($k>0$), $\et b_0=0$. The predicted crystal
$D(\lambda)=\bigoplus_{\mu}B(\mu)$, $\mu$ dominant with $\mu\le\lambda$, has one
copy of each such string. An \emph{isomorphism} $M(\lambda)\cong D(\lambda)$ is a
bijection $\psi$ preserving weights and satisfying $\psi(\ft x)=\ft\psi(x)$ and
$\psi(\et x)=\et\psi(x)$ (including $\ft x=0\iff\ft\psi(x)=0$, and the same for
$\et$).

\begin{lemma}\label{lem:hw}
Let $\psi:V\to D(\lambda)$ be an isomorphism. Then $\psi$ restricts to a
weight-preserving bijection between the highest-weight vertices
($\et x=0$) of $V$ and those of $D(\lambda)$. The highest-weight vertices of
$D(\lambda)$ are the $b_0$ of the copies of $B(\mu)$, one of each weight
$\mu\in\{\lambda,\lambda-2,\dots\}$. In particular $V$ must contain exactly one
highest-weight vertex of weight $\mu$ for each such $\mu$, and no others.
\end{lemma}
\begin{proof}
$\et x=0\iff\et\psi(x)=0$, and $\psi$ preserves weights. In $B(\mu)$ only $b_0$
has $\et b=0$, and its weight is $\mu$.
\end{proof}

\section{The counterexample \texorpdfstring{$\lambda=2$}{lambda = 2}}
\begin{proposition}\label{prop:M2}
$M(2)=\{Y_0^2,\ Y_0Y_1^{-1},\ Y_1^{-2}\}$ and $M(2)\cong B(2)$.
\end{proposition}
\begin{proof}
Direct computation from the definitions:
\begin{center}
\begin{tabular}{lcccll}
\toprule
$M$ & $\wt$ & $\varphi$ & $\varepsilon$ & $\ft M$ & $\et M$\\
\midrule
$Y_0^2$ & $2$ & $2$ & $0$ & $A_0^{-1}Y_0^2=Y_0Y_1^{-1}$ ($n_f=0$) & $0$\\
$Y_0Y_1^{-1}$ & $0$ & $1$ & $1$ & $A_0^{-1}Y_0Y_1^{-1}=Y_1^{-2}$ ($n_f=0$) & $A_0Y_0Y_1^{-1}=Y_0^2$ ($n_e=0$)\\
$Y_1^{-2}$ & $-2$ & $0$ & $2$ & $0$ & $A_0Y_1^{-2}=Y_0Y_1^{-1}$ ($n_e=0$)\\
\bottomrule
\end{tabular}
\end{center}
For instance, for $M=Y_0Y_1^{-1}$ the prefix sums $\sum_{k\le n}y_k$ are $0$
($n<0$), $1$ ($n=0$), $0$ ($n\ge1$), so $\varphi=1$, $n_f=0$; the values
$-\sum_{k>n}y_k$ are $0$ ($n<0$), $1$ ($n=0$), $0$ ($n\ge1$), so
$\varepsilon=1$, $n_e=0$. The set is closed under $\et,\ft$ and contains $Y_0^2$,
and every element is reached from $Y_0^2$ by $\ft$; so it is $M(2)$. The map
$Y_0^{2-k}Y_1^{-k}\mapsto b_k$ is an isomorphism onto $B(2)$.
\end{proof}

\begin{theorem}\label{thm:main}
$M(2)\not\cong B(2)\oplus B(0)$. Hence Conjecture 00000003824 is false.
\end{theorem}
\begin{proof}
$B(0)=\{b_0\}$ is a highest-weight vertex of weight $0$. By
Lemma~\ref{lem:hw}, an isomorphism would need a monomial $x\in M(2)$ with
$\et x=0$ and $\wt x=0$. The only element of weight $0$ is $Y_0Y_1^{-1}$, and
$\et(Y_0Y_1^{-1})=Y_0^2\ne0$. (Alternatively: $|M(2)|=3\neq4=|B(2)\oplus B(0)|$.)
\end{proof}

\section{Every \texorpdfstring{$\lambda\ge2$}{lambda >= 2} fails}
\begin{theorem}\label{thm:all}
For every $\lambda\ge1$,
$M(\lambda)=\{Y_0^{\lambda-k}Y_1^{-k}:0\le k\le\lambda\}\cong B(\lambda)$.
Consequently, for every $\lambda\ge2$,
$M(\lambda)\not\cong\bigoplus_{\mu\le\lambda}B(\mu)$, because $B(\lambda-2)$ is
missing.
\end{theorem}
\begin{proof}
Let $M_k=Y_0^{a}Y_1^{b}$ with $a=\lambda-k$, $b=-k$. If $0<k<\lambda$
($a>0>b$), the prefix sums are $0,a,a+b$, so $\varphi=a$ and $n_f=0$, and
$\ft M_k=A_0^{-1}M_k=Y_0^{a-1}Y_1^{b-1}=M_{k+1}$; the values $-\sum_{k>n}$ are
$-(a+b)$, $-b$, $0$, so $\varepsilon=-b=k$ and $n_e=0$, and
$\et M_k=A_0M_k=M_{k-1}$. For $k=0$: $\varphi(Y_0^\lambda)=\lambda$,
$\ft Y_0^\lambda=M_1$, $\varepsilon=0$. For $k=\lambda$: $\varphi(Y_1^{-\lambda})=0$,
$\varepsilon=\lambda$, $\et Y_1^{-\lambda}=A_0Y_1^{-\lambda}=M_{\lambda-1}$. So the set
$\{M_0,\dots,M_\lambda\}$ is closed, it is a single $\ft$-string from $M_0$, and
$M_k\mapsto b_k$ is an isomorphism with $B(\lambda)$. Its only highest-weight
vertex is $M_0$, of weight $\lambda$. For $\lambda\ge2$, $\lambda-2$ is dominant and
$\le\lambda$, and by Lemma~\ref{lem:hw} the predicted crystal requires a
highest-weight vertex of weight $\lambda-2$.
\end{proof}
For $\lambda\in\{0,1\}$ the only dominant $\mu\le\lambda$ is $\lambda$ itself and the
claim holds. Theorem~\ref{thm:all} is the $\mathfrak{sl}_2$ case of Kashiwara's
theorem~\cite{kas} that the connected component of a highest-weight monomial of
dominant weight $\lambda$ is isomorphic to $B(\lambda)$, in every type.
So the conjecture fails whenever some dominant $\mu<\lambda$ exists; for
example, for $\mathfrak{sl}_3$ and the adjoint weight $\theta=\varpi_1+\varpi_2$,
$M(\theta)$ has $8$ elements (it is $B(\theta)$), while $B(\theta)\oplus B(0)$
has $9$ (checked by \texttt{verify.py}).

\begin{remark}\label{rem:conv}
Conventions. (i) Nakajima~\cite{nak} uses variables $Y_{i,aq^k}$ and
$A_{i,aq^k}$ with a parity condition; after re-indexing these become
Kashiwara's $Y_m$, $A_m$ for a suitable choice of $c$ (in rank one there is
no choice), as explained in~\cite{kas}. (ii) The mirrored convention
(sums over $k\ge n$ instead of $k\le n$) is obtained by $m\mapsto-m$ and gives an
isomorphic crystal. (iii) Starting from $Y_s^\lambda$ instead of $Y_0^\lambda$ is a
shift of indices. In all cases $M(\lambda)$ is a single string $B(\lambda)$;
\texttt{verify.py} checks (i)--(ii) for $\lambda\le8$.
\end{remark}

\section{Other readings of ``\texorpdfstring{$M(\lambda)$}{M(lambda)}''}
The definition in the conjecture is informal, so we also consider the two
other readings we could think of in which $M(\lambda)$ is built from monomials
and has several components.

\paragraph{(U) Union reading.} $M_U(\lambda)$ is the union of the subcrystals
generated by \emph{all} highest-weight monomials of weight $\lambda$ (that is,
$\varepsilon(M)=0$, $\wt M=\lambda$). This fails for \emph{every}
$\lambda\in\mathbb N$: $Y_0^\lambda$ and $Y_1^\lambda$ (for $\lambda\ge1$), resp.\ $1$ and
$Y_1^{-1}Y_2$ (for $\lambda=0$), are two distinct highest-weight vertices of
weight $\lambda$, so $B(\lambda)$ would occur at least twice, contradicting
``exactly once'' (Lemma~\ref{lem:hw}). (By Kashiwara's theorem every component
of $M_U(\lambda)$ is a copy of $B(\lambda)$, so for $\lambda\ge2$ also no
$B(\mu)$ with $\mu<\lambda$ occurs; \texttt{verify.py} confirms this on all
highest-weight monomials of weights $0,\dots,3$ supported on $Y_0,\dots,Y_4$
with exponents in $[-2,3]$, finding $190$, $286$, $381$ and $467$ of them.)
Any larger reading, such as all monomials of weight $\le\lambda$, contains
$M_U(\lambda)$ and fails for the same reason.

\paragraph{(KR) Kirillov--Reshetikhin reading.} One might take for $M(\lambda)$
the monomial crystal attached to the Kirillov--Reshetikhin module with highest
weight $\lambda$. In type $A_1^{(1)}$ the KR modules are irreducible as
$U_q(\mathfrak{sl}_2)$-modules, and their highest-weight monomial (for
$\lambda=2$, $Y_0Y_1$ in the present indexing, i.e.\ $Y_aY_{aq^2}$) generates
$\{Y_0Y_1,\ Y_0Y_2^{-1},\ Y_1^{-1}Y_2^{-1}\}\cong B(2)$. So this reading
coincides with the main one, and $B(0)$ is again missing (Lean:
\texttt{kr2\_component}, \texttt{kr2\_fails}).

\paragraph{(P) Tensor-power reading.} ``Generated by monomial variables'' could
mean: $M_P(\lambda)$ is the subcrystal generated by the products
$c_0c_1\cdots c_{\lambda-1}$ with $c_j\in\{Y_{2j},Y_{2j+1}^{-1}\}$, each factor
being a copy $\{Y_{2j},Y_{2j+1}^{-1}\}\cong B(1)$. This set has $2^\lambda$
elements and realizes $B(1)^{\otimes\lambda}$. For $\lambda=2$ it is
$B(2)\oplus B(0)$ (highest-weight vertices $Y_0Y_2$, $Y_1^{-1}Y_2$), so under this
reading the case $\lambda=2$ is \emph{true}. For $\lambda=3$ it has the three
highest-weight vertices $Y_0Y_2Y_4$ (weight $3$), $Y_0Y_3^{-1}Y_4$ and
$Y_1^{-1}Y_2Y_4$ (both of weight $1$), i.e. $M_P(3)\cong B(3)\oplus B(1)^{\oplus2}$:
$B(1)$ occurs twice, so ``each exactly once'' fails (Lemma~\ref{lem:hw}).
In general the multiplicities are the Clebsch--Gordan numbers
(\texttt{verify.py}, $\lambda\le5$).

\begin{corollary}
For $\mathfrak{sl}_2$ and $\lambda=3$ the conjecture fails under the main
reading, under (U) and under (P).
\end{corollary}

\section{Lean formalization}
The directory \texttt{lean4/} is a self-contained Lean~4.19.0 project (core
library only).
\begin{itemize}
\item \texttt{Mon}: a monomial as a list of pairs $(m,y_m)$, $y_m\ne0$, indices
strictly increasing. \texttt{mulY}, \texttt{mulA} multiply by $Y_n^d$ and
$A_n^{s}$; \texttt{wt}, \texttt{phi}, \texttt{eps} implement the formulas
above (prefix/suffix maxima, the empty prefix/suffix included);
\texttt{firstHit}, \texttt{lastHit} compute $n_f$, $n_e$; \texttt{fT}, \texttt{eT}
are $\ft,\et$ with \texttt{none} for $0$.
\item \texttt{Reach s x}: $x$ is obtained from $s$ by finitely many $\ft,\et$.
\texttt{Mmain l = Reach (Y0pow l)} is $M(\lambda)$. \texttt{closeN} computes
closures, and \texttt{mem\_closeN\_iff} proves that a closed computed closure is
exactly the generated subcrystal.
\item \texttt{Target l}, \texttt{wtB}, \texttt{fB}, \texttt{eB}: the crystal
$\bigoplus_{\mu}B(\mu)$ with vertices $(\mu,k)$, $\mu$ dominated by $\lambda$
(\texttt{Dominated}: $\mu\le\lambda$, $\lambda-\mu$ even), $k\le\mu$.
\texttt{CrystalIso} is a bijection preserving $\wt$, $\ft$, $\et$;
\texttt{Decomposes V l} is the existence of such an isomorphism onto
\texttt{Target l}; \texttt{Conjecture} is $\forall\lambda$,
\texttt{Decomposes (Mmain l) l}.
\item \texttt{M2\_eq}: $M(2)$ is exactly the three monomials.
\texttt{iso\_M2\_B2}: $M(2)\cong B(2)$.
\item \texttt{conjecture\_00000003824\_clause\_false} proves
\[\neg\,\texttt{Decomposes (Mmain 2) 2},\]
and \texttt{conjecture\_00000003824\_false} proves $\neg$\,\texttt{Conjecture}.
\item \texttt{main\_fails\_all}: $\neg$\texttt{Decomposes (Mmain l) l} for every
$\lambda\ge2$, via the symbolic computation of $\ft,\et$ on the string
(\texttt{fT\_s1}, \dots, \texttt{InStr\_hw}).
\item \texttt{main\_fails\_all'}: the variant \texttt{Target'} without the parity
condition fails for every $\lambda\ge1$; \texttt{kr2\_fails}: the KR reading at
$\lambda=2$.
\item \texttt{union\_fails} (every $\lambda$), \texttt{prod3\_fails} and
\texttt{all\_readings\_fail\_at\_3} for the readings (U), (P).
\item Non-vacuity: \texttt{main\_holds\_0}, \texttt{main\_holds\_1},
\texttt{prod\_holds\_2} construct explicit isomorphisms, so
\texttt{Decomposes} is satisfiable and agrees with the cases where the claim is true.
\texttt{sanity} checks canonical form, $\varphi-\varepsilon=\wt$ and
$\et\ft=\mathrm{id}$, $\ft\et=\mathrm{id}$ on all components used.
\end{itemize}
There is no \texttt{sorry}, \texttt{native\_decide} or added axiom.
The command \texttt{\#print axioms} shows only the standard axioms
\texttt{propext} and \texttt{Quot.sound}; in addition,
\texttt{Classical.choice} appears for the three theorems
\texttt{main\_fails\_all}, \texttt{main\_fails\_all'} and
\texttt{all\_readings\_fail\_at\_3}.

\paragraph{Scope.} The Lean development is for $\mathfrak{sl}_2$. It does not
formalize Kashiwara's general theorem; it computes the relevant components
directly. The crystal axiom $\ft M=M'\iff\et M'=M$ is checked on the components
used but not proved for all monomials; it is only needed to identify the
subcrystal generated by $Y_0^\lambda$ with the undirected connected component,
which is a standard fact~\cite{kas}. The $\mathfrak{sl}_3$ example and the
enumerations quoted for (U) and (P) are in \texttt{verify.py} only.

\section{Independent check}
\texttt{verify.py} (Python~3 standard library) implements the monomial crystal
for an arbitrary Cartan matrix with Kashiwara's constants $c_{ij}$. It computes
$\varphi$, $\varepsilon$, $n_f$, $n_e$ by brute force over all $n$ in a window
around the support, which is a different method from the scans in Lean. It
checks: $M(\lambda)$ is a single string $B(\lambda)$ for $\lambda\le8$ in both
conventions of Remark~\ref{rem:conv}(ii), and the conjecture holds exactly for
$\lambda\le1$; the union reading (U); the tensor-power reading (P) with
Clebsch--Gordan multiplicities for $\lambda\le5$; and $|M(\varpi_1+\varpi_2)|=8$ for
$\mathfrak{sl}_3$.

\section{Reproduction}
\begin{verbatim}
cd lean4 && lake build     # Lean 4.19.0, core only
python3 verify.py          # standard library only
pdflatex report.tex
\end{verbatim}

\begin{thebibliography}{9}
\bibitem{kas} M.~Kashiwara, Realizations of crystals, in: \emph{Combinatorial and
Geometric Representation Theory}, Contemp.\ Math.\ 325, AMS, 2003, 133--139.
\bibitem{nak} H.~Nakajima, $t$-analogs of $q$-characters of quantum affine
algebras of type $A_n$, $D_n$, in: \emph{Combinatorial and Geometric
Representation Theory}, Contemp.\ Math.\ 325, AMS, 2003, 141--160.
\bibitem{kks} S.-J.~Kang, J.-A.~Kim, D.-U.~Shin, Monomial realization of crystal
bases for special linear Lie algebras, \emph{J.~Algebra} 274 (2004), 629--642.
\end{thebibliography}
\end{document}
